% authority: science_draft
% claim_status: draft/control
\documentclass[11pt]{article}
\usepackage[margin=1in]{geometry}
\usepackage{amsmath,amssymb,booktabs,array,enumitem,hyperref}
\hypersetup{hidelinks}
\setlist{nosep}

\title{V22 P1-T04 Gate-B-Only Physics Lock\\
\large Dependency closure, physical typing, downstream parking, and historical preservation}
\author{Smuggling Auditor parent with Process Integrity Auditor perspective}
\date{2026-08-09}

\newcommand{\Eligible}{\operatorname{Eligible}}
\newcommand{\Positive}{\mathsf{POSITIVE}}
\newcommand{\Parked}{\mathsf{PARKED}}
\newcommand{\Typed}{\operatorname{PhysTyped}}
\newcommand{\Preserve}{\operatorname{Preserve}}

\begin{document}
\sloppy
\maketitle

\begin{center}
\fbox{\parbox{0.92\textwidth}{
\textbf{Status and authority.} This is the draft/control V22 P1-T04
admission policy. It makes Gate B the sole active physics gate and parks new
downstream execution. It is not a positive Gate B verdict, a new Gate C
review, an effective metric, an Einstein equation, benchmark evidence,
external review, ontology adoption, or a completed derivation.
}}
\end{center}

\section{Decision and live burden}

The active derivation milestone is \texttt{effective\_metric\_g\_eff}. Its
direct burden is the source-derived metric law: construct or constrain a
target-atlas-free source-to-geometry map that meets all eight Gate B criteria.
No current record supplies a genuinely positive Gate B result. The historical
V21 separating certificate records zero of eight Gate B criteria met and keeps
same-milestone continuation open. Therefore Gate B is the sole active physics
gate. A task is not admissible merely because it is interesting, inexpensive,
already implemented, protected by a human decision, or accepted by a
validator.

The only admissible physics task classes are:
\begin{enumerate}
  \item proving a Gate B necessary condition;
  \item constructing a source-derived Gate B candidate;
  \item proving a scoped Gate B obstruction;
  \item comparing materially distinct minimal source extensions; and
  \item testing a Gate B candidate's robustness or uniqueness.
\end{enumerate}
Every admitted task must name its direct Gate B burden, assumptions, evidence
type, failure result, and the absence of protected or process-derived truth
promotion.

\section{Gate-lock dependency theorem}

Let $R$ be the acyclic route graph and let $B$ denote Gate B. For a route $r$,
let $\operatorname{Pred}(r)$ be its required positive scientific predecessor
records. Let $\mathcal A_B$ be the five Gate-B-direct classes above, and let
$D_B$ be the routes downstream of $B$: new Gate C, Gate D, Gate E,
Einstein-route physics, benchmark execution, and their conditional P5/P6
packages. Define
\[
 \Eligible(r)=
 \begin{cases}
  1,&r\in\mathcal A_B\text{ and its direct Gate B burden is typed};\\
  1,&r\notin D_B\text{ and every required predecessor is positive};\\
  0,&\text{otherwise}.
 \end{cases}
\]
For $r\in D_B$, positive Gate B evidence is a necessary predecessor. Gate D
also requires a genuinely positive Gate B decision; Gate E and benchmark
execution require the later positive Gate D path. A protocol may be defined
early only as control text: this does not authorize its physics execution.

\paragraph{Gate-lock dependency-closure theorem.}
If there is no genuinely positive Gate B evidence record, then
$\Eligible(r)=0$ for every $r\in D_B$. A protected decision, postulate,
target-GR input, familiar label, validator PASS, checkpoint, or internal review
cannot change this result.

\paragraph{Proof.}
Every $r\in D_B$ has Gate B itself or a route whose predecessor closure
contains Gate B. The definition of eligibility requires all scientific
predecessors to be positive. Because Gate B is not positive, the conjunction
is false. Induction along a topological ordering of the acyclic route graph
propagates the false predecessor to Gate D, Gate E, Einstein, benchmark, and
conditional execution descendants. The no-postulate policy makes protected
and operational records evidence-invariant, so none supplies the missing
scientific predecessor. Hence every downstream route remains parked.
$\square$

This theorem is a dependency-policy result, not a proof that Gate B can never
be passed. It leaves same-milestone candidate, obstruction, extension, and
robustness research open.

\section{Exact downstream parking}

New Gate C execution is parked; the existing protected P7 Gate C record is
preserved at its exact finite source-matter scope and is not rerun. Gate D,
Gate E, Einstein-route physics, and benchmark execution are parked. All P5
conditional packages remain ineligible until P4-T05 records genuinely positive
Gate B evidence. P6-T02 remains ineligible until P5-T03 records genuinely
positive Gate D evidence. P6-T01 may define a protocol early only as a
nonexecuting control artifact. P6-T03 remains parked until its declared
positive or terminal-failure predecessor path actually exists. No conditional
P5/P6 physics package executes early.

Candidate processing is also fail closed. A candidate stops when its
gauge-reduced principal polynomial is absent or nonhyperbolic, its cone is
unstable, propagation is non-universal, a hidden metric-equivalent primitive
is detected, or its continuum/refinement control is absent. At most one source
extension family may be active. A fallback needs a recorded termination and a
new selector decision.

\section{Physical-type gate}

For a physical label $\ell$ applied to an object $x$, define
\[
 \Typed(\ell,x)=T_s(x)\wedge O_\ell(x)\wedge W_\ell(x)
 \wedge U_\ell(x)\wedge D_\ell(x)\wedge F_\ell(x),
\]
where $T_s$ is an explicit source type; $O_\ell$ gives operational semantics;
$W_\ell$ is a source-to-world or source-to-observable map; $U_\ell$ gives
units, scale, and calibration; $D_\ell$ gives domain, regime, and equivalence
scope; and $F_\ell$ gives a falsifiable validation or discrimination rule.
Label-specific obligations are additional: a clock needs repeatable rate and
comparison semantics; a signal needs preparation, propagation, detection, and
causal semantics; free fall needs a universal trajectory or response law;
gravity needs a universal interaction or effective-geometry relation; a metric
needs a tensorial nondegenerate bilinear structure with transformation and
gluing laws; and a compact source needs physical matter, localization,
exterior matching, and controlled compactness.

\paragraph{Finite-token physical-typing lemma.}
A finite token, vertex name, graph label, register value, or renamed algebraic
object does not by itself satisfy $\Typed(\ell,x)$ for \texttt{clock},
\texttt{signal}, \texttt{free-fall}, \texttt{gravity}, \texttt{metric}, or
\texttt{compact source}.

\paragraph{Proof.}
A token supplies at most a syntactic name and a finite carrier. The definition
requires six independently declared witnesses, including operational semantics
and a source-to-world map. Naming alone constructs none of those witnesses.
Therefore at least one required conjunct is absent, so the physical-type
predicate is false. The claim is contract-relative: a later source-derived
construction may supply the witnesses and be tested anew. $\square$

\section{Gate B criteria and direct burden}

Every candidate path remains answerable against the same eight criteria:
target-atlas-free construction; a gauge- and constraint-reduced hyperbolic
principal polynomial; one stable Lorentzian cone or physically derived
equivalence class; universal compatibility across declared matter sectors;
operational clocks, rods, and scale; tensorial transformation and
local-to-global gluing; finite-variation and coarse-graining robustness; and no
hidden metric, tetrad, coframe, constitutive tensor, quadratic form, principal
polynomial, or equivalent target primitive. A Director can explain an eligible
task only by pointing to one of these burdens or to a necessary condition,
scoped obstruction, minimal extension comparison, or robustness/uniqueness
test that directly narrows them.

\section{Historical P8/P9 preservation theorem}

Let $M_{21}$ be the immutable V21 baseline manifest at commit
\path{233e5dd7024fc068032d0afe86d85dc25e2246e9} and tree
\path{a7d9c9448de8e204643b093878ba4d84bd58f020}. Let $H_{89}$ be the
66 manifest entries under the sixteen historical P8/P9 task roots registered
by this packet. Define $\Preserve(H_{89})$ to require every live path to exist
with exactly its baseline SHA-256.

\paragraph{Historical projection and nonpromotion theorem.}
If $\Preserve(H_{89})$ holds, the P8/P9 archive remains complete historical
negative/control evidence, but its preservation cannot create a positive Gate
record or make a parked route eligible.

\paragraph{Proof.}
The baseline manifest fixes path and byte hash, so preservation is an identity
map on the archived records. Those records state Gate D NOT READY; six P9
benchmark cases INCONCLUSIVE; zero benchmark passes; zero qualifying
independent replications; and Gate E NOT READY. The identity map cannot change
any status field. The eligibility function reads positive predecessor evidence,
not mere record existence. Therefore archive preservation neither reruns the
historical packages nor promotes them into active or positive evidence.
$\square$

The task-local validator checks all 66 live paths and hashes, not only the
representative registered sources. P8/P9 outputs remain readable for negative
results, control comparisons, provenance, and future materially distinct
reopening analysis.

\section{Disposition}

P1-T04 is complete at draft/control scope when the gate-lock policy, admission
matrix, parked-route registry, physical-type checklist, historical index,
fixtures, independent internal process-integrity audit, registrations, and
validation pass. The active routing milestone changes to
\texttt{effective\_metric\_g\_eff}; the Distance-to-GR ledger does not change.
The next package may be P2-T01, a separately admitted Gate-B-direct necessary-
condition theorem. Selection is not execution. No ontology, metric, coupling,
Einstein, Gate, benchmark, proof, review, publication, release, push, or
completed-derivation authority is created.

\end{document}
